\chapter{引言}
\begingroup
\justifying
\setlength{\parindent}{0pt}
\setstretch{1.5}
\setlength{\parskip}{0.5\baselineskip}
\titlespacing{\chapter}{0pt}{0pt}{0pt}
\titlespacing{\section}{0pt}{0pt}{0pt}

电子投票承诺带来更快的计票速度、更广泛的可达性以及可审计的记录，这些优势使其多年来持续成为选举改革议程中的关键方向。然而，真正实现上述目标远比表面承诺更困难。一个可信系统必须同时保证选票保密性、强制执行正确的状态迁移、向外部公开足够的可验证信息，并且对管理员与选民都具备可用性。这些要求彼此牵制，而多数现实部署恰恰在这种张力下暴露问题。

中心化电子投票平台在可用性与处理速度方面通常表现良好，但其代价是将信任负担转移给后端运营方。运营方控制日志、执行状态迁移并发布计票结果，而外部独立验证的边界取决于其愿意公开的信息范围。当结果出现争议时，调查通常会停留在特权访问边界之外，而非通过公开可验证机制完成裁决。

本论文采用不同路径：将可信性从“对运营方的制度性信任”转向“链上执行约束与密码学证明”的可核验机制；其中计票密钥托管仍是显式信任前提。系统将 ElGamal 同态加密、Groth16 零知识证明和基于承诺的 Merkle 包含性证明整合为单一端到端流程。目标并非理论性草图，而是可编译、可部署、可测试并可独立审阅的工作系统。

\section{问题陈述}

构建一个最小化对受信运营方依赖的投票系统，需要同时满足一组方向并不完全一致的要求：
\begin{itemize}
    \item 在提交、存储与计票全过程保障选票机密性；
    \item 在不泄露选票内容的前提下，对投票构造与计票声明执行正确性校验；
    \item 以可篡改留痕方式强制执行选举状态迁移；
    \item 提供选民验证路径，使其可由本地回执结合公开审计包完成包含性验证，并将审计包根与链上 \texttt{merkleRoot} 比对；验证步骤不依赖地址检索。
\end{itemize}

设 $V$ 为已投选票数、$N$ 为候选人数。每一张被接受的选票都必须编码且仅编码一个有效选择，聚合密文必须导出一致计票结果，且公开结果必须附带可由任何人独立核验的密码学证据。

\section{研究动机}

若干现实技术进展使上述组合方案具备可行性与研究价值。区块链账本以确定性执行维护仅追加事件日志，可消除中心化后端长期暴露的一类攻击面，即选举状态的静默重写 \cite{nakamoto2008,wood2014}。现代 zk-SNARK 系统已达到证明体积小、链上验证可纳入智能合约预算的成熟阶段 \cite{groth2016}。面向电路友好的哈希函数与 Merkle 结构进一步提供可扩展的选后包含性验证能力，而无需昂贵链上存储 \cite{poseidon2019,merkle1988}。

综合来看，这些工具打开了一条新的设计路径：在不牺牲验证能力的条件下更好地保护隐私，并将系统可信性从机构承诺转向密码学校验与公共执行轨迹。

\section{研究目标}

本研究聚焦于五项结果：
\begin{itemize}
    \item 构建合约级状态机，对选举生命周期迁移进行显式约束；
    \item 基于 BabyJubJub 上的一位有效 ElGamal 加密实现隐私投票提交；
    \item 通过 Groth16 在链上验证投票合法性与计票正确性；
    \item 建立“回执驱动”的选票包含性验证路径：选民使用本地回执与公开审计包完成核验，并将审计包根与链上 \texttt{merkleRoot} 对比；验证步骤无需额外身份暴露；
    \item 对实现效果进行诚实评估，明确指出已完成能力与尚存缺口。
\end{itemize}

\section{研究问题}

系统的实现与评估围绕以下三个问题展开：
\begin{itemize}
    \item \textbf{RQ1}：在不依赖受信中间件的前提下，将加密 one-hot 选票与基于零知识证明的校验整合到浏览器到链流程中是否可行？
    \item \textbf{RQ2}：相较于中心化日志审计，该流程是否提供更强的可验证性（判据包括：无需特权数据即可独立重放计票声明；选民无需服务器交互即可完成包含性验证）？
    \item \textbf{RQ3}：在目标安全保证中，哪些在当前代码中已完整实现，哪些仍受具体工程缺口限制？
\end{itemize}

表 \ref{tab:rq-mapping} 给出了研究问题与验证位置的对应关系。

\begin{table}[htbp]
\centering
\small
\caption{研究问题与验证证据映射}
\label{tab:rq-mapping}
\begin{tabular}{L{2.0cm} L{5.0cm} L{5.8cm}}
\toprule
RQ & 评估重点 & 主要证据位置 \\
\midrule
RQ1 & 客户端密码学、证明生成与合约验证的端到端集成可行性 & 第 4 章实现流程与第 5 章场景验证 \\
RQ2 & 相较中心化电子投票的可验证性增益（按可独立重放性与无服务器验证能力评估） & 第 2 章定位（表 \ref{tab:approach-comparison}）与第 5 章安全分析 \\
RQ3 & 功能完成与高保障完成之间的差距分析 & 第 5 章未完成缺口与附录~\ref{tab:acceptance-checklist} 清单 \\
\bottomrule
\end{tabular}
\end{table}

\section{研究范围与边界}

本系统面向单场选举部署模型，并采用管理员控制的计票密钥。该原型定位为可复现的研究演示系统，而非面向国家级选举的生产平台。

以下内容不在本论文实现范围：抗胁迫机制、门限解密仪式、选举管理去中心化治理、合约字节码形式化验证。上述方向均作为后续工作而非已完成能力。

\section{主要贡献}

本研究产出五项具体成果：
\begin{itemize}
    \item 提供可跨栈复现实验原型（Solidity 合约、Circom 电路、浏览器端 JavaScript 证明、Python 密码学验证），并支持编译、部署与独立测试；
    \item 实现投票提交机制，将 Poseidon 承诺与密文哈希作为公开证明信号进行绑定；
    \item 构建计票流程：链下同态聚合与 ZKP3 验证完成后，发布 Merkle 根以支持承诺集合可审计性（在第 5 章分析的实现假设下成立）；
    \item 给出选后审计路径，使选民可基于本地回执与导出数据包验证选票包含性，验证流程不依赖地址检索；
    \item 提供需求驱动的缺口评估，对未解决风险进行显式命名，而非将其笼统并入一般性局限。
\end{itemize}

\section{论文结构}

第 2 章综述相关工作并构建协议依赖的密码学基础。第 3 章定义系统模型、协议流程与信任假设。第 4 章说明实现决策与模块级行为。第 5 章报告评估结果，包括性能测量、需求可追踪性与风险分析。第 6 章给出结论并提出分优先级的加固路线。

\section{本章小结}

本章明确了问题背景，提出三项研究问题，并界定了安全结论的解释边界。同时，本章给出了主要贡献项，后续章节将基于实现与评估证据逐一验证。第 2 章将把本研究置于更广泛的电子投票与密码学文献语境中，并进一步展开协议设计所依赖的理论基础，尤其是 Groth16。

\endgroup
